\section{2025 应用窗口：为什么 Lovart 会在此时出圈}

Lovart 的出圈不是孤立事件，它出现在 2025 年 Agent 应用窗口打开之后。底座模型已经足够强，用户也开始理解 Agent 的价值，但通用 Agent 还没有覆盖所有专业工作流。这个窗口给垂直应用留下机会：如果能在专业场景里做出明显好于通用工具的体验，就可能快速获得全球关注。

\begin{knowledgebox}{应用窗口的三个条件}
\begin{tabularx}{\textwidth}{p{0.20\textwidth}p{0.34\textwidth}X}
\toprule
条件 & 解释 & Lovart 对应点 \\
\midrule
模型能力够用 & 底座模型能理解需求并生成高质量内容 & 设计生成和多轮修改可用。 \\
通用入口未吃完 & ChatGPT/Manus 还不能深做所有垂直场景 & 设计师工作流仍有空位。 \\
用户心智成熟 & 用户愿意尝试把专业任务交给 AI & 设计师和创作者开始寻找 AI 工作台。 \\
\bottomrule
\end{tabularx}
\end{knowledgebox}

\begin{warningbox}{窗口期也意味着倒计时}
当一个垂直 Agent 跑出来，通用 Agent、模型公司和竞品都会学习它。窗口打开的同时，窗口也在缩短。Lovart 必须在被学习之前沉淀工作流和心智。
\end{warningbox}

\subsection{本章小结}

Lovart 的出圈来自模型能力、用户心智和垂直场景空位的叠加。它的机会不是永久的，必须尽快从窗口红利变成工作流壁垒。

